\section{Preliminaries}
This section explains the models and methods the project is built on.

\subsection{Transformer Encoders and Model Distillation}
The transformer \cite{vaswani2017attention} is a type of neural network that reads text with self-attention instead of word by word. BERT \cite{devlin2019bert} showed a simple recipe that this project relies on: first train a model on a very large amount of general text, then fine-tune it for a specific task on a much smaller amount of labelled data. The recipe suits us, because 510 contracts is a small dataset for deep learning.

Because all our work runs on one laptop, we use DistilBERT \cite{sanh2019distilbert}, a smaller, six-layer version of BERT that keeps about 97\% of BERT's language ability while being 40\% smaller and 60\% faster. The published CUAD results use much larger models, so we give up some accuracy in exchange for running everything on normal hardware. We point out this size difference whenever we discuss our results.

\subsection{Long-Document Strategies and Multiple-Instance Learning}
\label{sec:lit_long}
Contracts are much longer than a transformer can read at once. There are three usual ways to deal with this: (i) \emph{retrieval}, where a quick, cheap method picks the few parts of the text most likely to contain the answer and only those go to the transformer; (ii) \emph{sliding windows}, where the model reads every part of the text and the scores are combined into one answer for the whole document; and (iii) \emph{sparse-attention models} such as Longformer, which can read a longer piece of text at once.

The third option is often dismissed in one line as ``too expensive'', which is not enough of a reason. We did not run a Longformer experiment, so what follows is a design choice. The one part of it we could check cheaply on our own laptop was the cost, and Section~\ref{sec:lfbench} reports that check.

We had three reasons, and only the third is about cost. The first and most important is that Longformer would not solve the problem. It can read 4{,}096 tokens at once. At about four characters per token, our typical contract of 33{,}143 characters is about 8{,}300 tokens and the longest is about 85{,}000, so Longformer would still see only half of a typical contract and one-twentieth of the longest one. It makes the window eight times wider, which is still not wide enough for this data. We would still need windows and a way to combine their scores, which is the design we ended up using anyway, only with a much more expensive model underneath. This reason comes from arithmetic about the dataset and not from an experiment, but the arithmetic is clear: CUAD contracts are about ten times too long for a wider window to fix.

The second reason is availability. Our whole system uses small (distilled) models so that it can be trained and run on one laptop. We do not know of a small version of Longformer that is as reliable and well supported as DistilBERT, so using Longformer would have meant giving up the small models that make the rest of the project easy to repeat.

The third reason is cost in practice. We believed that fine-tuning a Longformer with 148.7M parameters at 4{,}096 tokens was not possible on one Apple-Silicon laptop. We were also worried about special attention code on a non-CUDA machine, after DeBERTa-v3 produced NaN (broken) values on the MPS backend (Section~\ref{sec:engineering}). When we later measured it (Section~\ref{sec:lfbench}), the memory part of this belief turned out to be wrong, since one training step does fit at batch size 1. The time part was right by a very large margin. We kept our original reasoning in the text and reported the correction next to the measurement, so the reader can see both. The app adds its own cost: \emph{Clausify} scores every window of an uploaded contract for all 41 categories while the user waits, so it is only usable if this takes seconds. That can be measured without any training, and Section~\ref{sec:lfbench} gives the speed and peak memory of Longformer-base at 4{,}096 tokens on our laptop, next to the DistilBERT setup we used.

On a CUDA computer with more memory, the third reason would be weaker and the second might disappear, but the first would still hold. The choice is backed by one measurement and some arithmetic. It does not prove that sparse attention cannot work on CUAD. A properly trained Longformer with windows would be a fair comparison to ask for, and we list it as future work.

Reading a document in windows and taking the highest window score is known as \emph{multiple-instance learning} (MIL). The document is treated as a bag of windows, and the bag counts as positive if any one window is positive, which is what taking the maximum score does. MIL has a well-known weakness, and our experiments confirmed it: it produces too many false alarms. A document has dozens of windows, so the chance that some window fires by mistake grows with the length of the document, and precision drops.
\subsection{Text-to-Text Models for Summarization}
T5 \cite{raffel2020t5} treats every language task as ``text in, text out'', which suits clause simplification: legal text goes in and a simple sentence comes out. FLAN-T5 \cite{chung2022flan} is a version of T5 trained on many tasks with instructions, so it already gives useful answers to a simple prompt such as ``summarize in plain English:'' without extra training. We score summaries with ROUGE-L \cite{lin2004rouge}, which measures how many words the output shares, in the same order, with a reference sentence. Chapter~5 shows why, when there are only a few reference sentences, this score has to be checked by reading the outputs.

\section{Review of Existing Research}
This section reviews earlier work on legal language processing, legal risk, plain-language legal writing, and the accuracy of machine-written summaries.

\subsection{Legal NLP and the CUAD Benchmark}
People have long wanted to use language technology on legal text, and progress was slow for a simple reason: labelling legal text needs lawyers, and lawyers are expensive. CUAD \cite{hendrycks2021cuad}, released by The Atticus Project \cite{atticus2021} at the NeurIPS 2021 Datasets and Benchmarks track, changed this. It provides 510 real business contracts from public SEC filings, in which experienced lawyers marked every piece of text that belongs to one of 41 important clause categories, such as governing law, non-compete, limits on liability and IP assignment. The data uses the SQuAD format \cite{rajpurkar2016squad}: each (contract, category) pair is asked as a question about the contract text, and the marked clauses are the answers, with their exact character positions.

The original CUAD paper treats the task as question answering and finds that bigger models do much better. Its best results come from the extra-large DeBERTa models, and even those scores are modest, which shows how hard it is to find legal clauses. Two features of CUAD affect any model built on it. Categories are very uneven: some appear in almost every contract, while others appear only about a dozen times in the whole dataset. Contracts are also long, far longer than a transformer can read at once, so every practical system first has to work out how a model that reads 512 tokens at a time can read a 50{,}000-character contract.
\subsection{Legal Risk Assessment and Decision Support}
\label{sec:lit_risk}
The risk part of this project has the least earlier work behind it, which is a weakness of our work. Since we tell a non-lawyer which clauses to worry about, the obvious question is what this advice is based on. We did not use an existing framework. We wrote the risk table ourselves (Section~\ref{sec:riskcriteria}), which is why we call it a prototype.

Earlier work does explain why this is hard. Legal language benchmarks usually treat ``risk'' indirectly, by sorting text into legally meaningful categories without judging how serious something is. LexGLUE \cite{chalkidis2022lexglue} collects tasks such as predicting court decisions and classifying laws, and LEGAL-BERT \cite{chalkidis2020legalbert} shows that models pre-trained on legal text do better on them. None of these tasks asks how dangerous a clause is for a particular person signing it. CUAD itself \cite{hendrycks2021cuad} describes its categories as the clauses lawyers look for when checking a deal, which marks them as important without saying how important.

How serious a clause is remains an open question, and it is hard to learn from data alone, for three reasons. Risk depends on which side you are on: the same unlimited indemnity is dangerous for the party giving it and useful for the party receiving it, so a risk label only makes sense once you pick a side, a problem Section~\ref{sec:weaker} handles only by making an assumption. Risk depends on the situation, since a \$50{,}000 limit on liability means little to a large company but can ruin a one-person business. And risk depends on the country, since a non-compete that is binding in one country may be invalid in another. Our table gives one fixed level per category, so it ignores all three on purpose.

This affects how our work should be read. Our risk table should be compared with expert opinion, because we know of no published model that covers CUAD's 41 categories, and that expert check has not been done yet (Section~\ref{sec:taxvalid}). Work on explainable legal AI argues that tools in this area should show their reasoning instead of just giving a score. Our table does this: every level comes with a written reason that can be questioned line by line (Appendix~\ref{app:taxonomy}). That openness makes the table easy to check, though it does not make the table correct.
\subsection{Plain-Language Legal Communication}
\label{sec:lit_plain}
Chapter~5 finds that our summarizer does not really summarize, and the closest earlier work is about simplifying text. The two goals differ. Summarizing makes text shorter, while simplifying says the same thing in easier words. For a contract clause, shorter is often the wrong goal, because the conditions and exceptions a summary would drop are often what decides what the clause does.

The closest earlier work is Manor and Li \cite{manor2019plain}, who collected plain-English summaries of contract sections and found that models trained on them tend to give general, repeated answers, an early version of the problem our template test measures. Measuring how easy a text is to read is a much older field \cite{flesch1948}, and it offers a check our evaluation lacks: that the output is easier to read, and not only shorter. We do not report any readability score. Looking back, that is a gap, since this part of the project exists to make clauses easier to understand.

Our results agree with the main lesson of this work: scores based on reference sentences measure the reference sentences. With only 41 different target sentences, ROUGE-L rewards the model for recognising the category. Section~\ref{sec:summ} measures this directly by changing the reference sentences while keeping everything else the same.
\subsection{Faithfulness and Hallucination in Generated Summaries}
\label{sec:lit_faith}
When a system shows machine-written sentences next to a user's own contract, whether those sentences are true is a safety issue. Research on summarization has shown that fluent machine-written summaries often say things that are not in the source. Maynez et al.\ \cite{maynez2020faithfulness} find that many generated summaries contain content the input does not support, and that ROUGE agrees poorly with human judgements of accuracy. Kry\'{s}ci\'{n}ski et al.\ \cite{kryscinski2020factcc} suggest checking factual accuracy as a separate task. Scores such as BERTScore \cite{zhang2020bertscore} measure meaning better than simple word overlap, but they do not check facts either.

Our system has an unusual version of this problem. A model that only outputs fixed template sentences cannot invent details in the usual way, because it can only produce one of 41 sentences written by people, and each one correctly describes its category. It fails in the opposite way. It says nothing about the specific clause, and when it picks the wrong template (30\% of the time, Section~\ref{sec:summ}) the output is fluent and confident but describes a completely different kind of clause. For a legal assistant this may be worse than an invented detail, because the output gives no sign of the error. The sentence is well written and true of its own category, and only a comparison with the clause shows that it does not describe it.

Our evaluation is therefore missing the checks this research recommends. We report ROUGE-L against two sets of reference sentences and we read outputs by hand. We do not report BERTScore, an automatic fact check, or a human judgement of whether the legal meaning is kept, and for a part that non-lawyers will see, the human judgement matters most. Section~\ref{sec:futurework} lists it as future work.

\section{Summary of Key Findings}
\label{sec:litgap}
Earlier CUAD work stops once the clause is found. A person signing a contract probably needs more: which clauses matter, how much, and what they mean in plain language. This project tries to add that layer, with a prototype risk table for each category, plain-English template sentences and a demo app. It also re-tests, on a strict test set, the common belief that a transformer must beat simple word-based methods at deciding which clauses a contract contains. Of the three added parts, the results support only the demo app as working. No expert has checked the risk table, and the plain-English part outputs category templates instead of summaries, so neither is solved. The added layer mostly gives a measurement of how far it is from working.

Table~\ref{tab:related} compares this work with the earlier work above. The first difference is scope. Earlier work stops at finding the clause text, while our output is a risk-sorted, plain-English report, which is the part a non-lawyer can act on. The second is in the ``Long inputs'' column. Every system on CUAD must choose how to handle long documents, and usually the choice is just stated. We measured ours before any training: the 64.0\% retrieval limit in Section~\ref{sec:ceiling} ruled out retrieval and led us to multiple-instance learning. The Model column also shows a cost. Our small models (67M parameters) are about ten times smaller than the DeBERTa-xlarge model behind CUAD's best published result, so our extraction scores cannot be compared directly with theirs, and we do not claim they can.

\begin{table}[!ht]
\centering
\caption{Related work}
\label{tab:related}
\footnotesize
\setlength{\tabcolsep}{3.5pt}
\renewcommand{\arraystretch}{1.25}
\hyphenpenalty=10000
\newcolumntype{L}[1]{>{\raggedright\arraybackslash}p{#1}}
\newcommand{\ds}[1]{\newline{\scriptsize\color{gray!70!black}#1}}
\begin{adjustbox}{max width=\textwidth}
\begin{tabular}{@{}L{25mm} L{23mm} L{21mm} L{25mm} L{17mm} L{16mm} L{23mm}@{}}
\toprule
\textbf{Work \textnormal{(data)}} & \textbf{Task} & \textbf{Model} & \textbf{Long inputs} & \textbf{Risk} & \textbf{Plain English} & \textbf{Evaluation}\\
\midrule
Hendrycks et al.\ \cite{hendrycks2021cuad}\ds{CUAD, 510 contracts}
 & Span extraction & DeBERTa-xlarge & Full-document QA with a large model & No & No & Held-out; precision at 80\% recall\\
\addlinespace
Chalkidis et al.\ \cite{chalkidis2020legalbert}\ds{several legal corpora}
 & Domain pre-training and classification & LEGAL-BERT & Truncation to 512 tokens & No & No & Per-task held-out\\
\addlinespace
Chalkidis et al.\ \cite{chalkidis2022lexglue}\ds{LexGLUE, 7 tasks}
 & Legal language understanding & BERT and RoBERTa models & Depends on the task & Indirect (task labels) & No & Multi-task benchmark\\
\addlinespace
Manor and Li \cite{manor2019plain}\ds{TL;DR-Legal style}
 & Plain-English summarization & Seq2seq & Section-level inputs & No & \textbf{Yes} & ROUGE and manual review\\
\addlinespace
Beltagy et al.\ \cite{beltagy2020longformer}\ds{long-document benchmarks}
 & New architecture & Longformer & Sparse attention, 4{,}096 tokens & No & No & Task benchmarks\\
\addlinespace
Retrieval--QA pipelines\ds{various}
 & Span extraction & Various & Retrieval then reader (ceiling usually assumed) & No & No & Various\\
\midrule
\rowcolor{blue!6}
\textbf{This project}\ds{CUAD, 510 contracts}
 & Presence, span, category risk and template text
 & DistilBERT and FLAN-T5 (small)
 & Sliding-window MIL; ceiling \emph{measured}, sparse attention \emph{benchmarked}
 & \textbf{Yes} (category level, prototype)
 & \textbf{Yes} (template based)
 & Contract-level held-out; cluster bootstrap; seed, budget and split controls; end-to-end\\
\bottomrule
\end{tabular}
\end{adjustbox}

\vspace{4pt}
{\footnotesize\singlespacing \emph{Note.} ``Long inputs'' is how each work handles documents longer than the encoder window; ``Evaluation'' is the level at which results are reported.\par}
\end{table}

